\documentclass[11pt,a4paper]{article}
\usepackage[utf8]{inputenc}
\usepackage[T1]{fontenc}
\usepackage[catalan]{babel}
\usepackage{lmodern}
\usepackage[sfdefault,lining]{FiraSans}
\usepackage[a4paper,top=2.3cm,bottom=2.2cm,left=1.8cm,right=1.8cm,headheight=1.3cm,headsep=0.4cm,footskip=1cm]{geometry}
\usepackage{amsmath,amssymb}
\usepackage{icomma}
\usepackage{graphicx}
\usepackage[table]{xcolor}
\usepackage{tikz}
\usepackage{fontawesome5}
\usepackage{ccicons}
\usepackage{fancyhdr}
\usepackage{lastpage}
\usepackage{tasks}
\usepackage{needspace}
\usepackage[most]{tcolorbox}
\usepackage{array}
\usepackage[hidelinks,pdftitle={Fitxa d'exercicis de la unitat 7: Funcions elementals. 4t d'ESO Matemàtiques A},pdfsubject={Matemàtiques A 4t ESO},pdfkeywords={funcions lineals, proporcionalitat directa, funció afí, funció constant, pendent, equació punt-pendent, funció quadràtica, paràbola, vèrtex, funcions radicals, proporcionalitat inversa, hipèrbola, asímptotes, translacions, funció exponencial, problemes amb funcions}]{hyperref}

\definecolor{navy}{HTML}{1B2748}
\definecolor{accent}{HTML}{E8603A}
\definecolor{teal}{HTML}{1F9E90}
\definecolor{soft}{HTML}{EEF1F8}
\definecolor{muted}{HTML}{5D6685}

\setlength{\parindent}{0pt}
\setlength{\parskip}{3pt}
\settasks{label=\alph*), label-format=\bfseries\color{navy}, label-width=1.3em, item-indent=1.8em, column-sep=0.6em, after-item-skip=5pt, before-skip=3pt, after-skip=4pt}

% ---------- capçalera i peu ----------
\newcommand{\llicencia}{\href{https://creativecommons.org/licenses/by-nc-sa/4.0/deed.ca}{\ccbyncsa}\ \ Llicència Creative Commons Reconeixement-NoComercial-CompartirIgual 4.0 Internacional (CC BY-NC-SA 4.0)}
\pagestyle{fancy}
\fancyhf{}
\fancyhead[L]{\raisebox{-0.35\height}{\includegraphics[height=0.95cm]{logo-gva.png}}}
\fancyhead[C]{\small\color{muted}Fitxa d'exercicis · Unitat 7: Funcions elementals · 4t ESO A}
\fancyhead[R]{\raisebox{-0.35\height}{\includegraphics[height=0.95cm]{logo-vermellar.png}}}
\fancyfoot[L]{\scriptsize\color{muted}\llicencia}
\fancyfoot[R]{\small\color{muted}\thepage}
\renewcommand{\headrulewidth}{0.4pt}
\renewcommand{\headrule}{\color{muted!40}\hrule height \headrulewidth}
\fancypagestyle{primera}{\fancyhf{}\renewcommand{\headrulewidth}{0pt}
  \fancyfoot[L]{\scriptsize\color{muted}\llicencia}
  \fancyfoot[R]{\small\color{muted}\thepage}}

% ---------- nivells i exercicis ----------
\newcommand{\estrelles}[1]{{\color{accent}\ifcase#1\or\faStar\faStar[regular]\faStar[regular]\or\faStar\faStar\faStar[regular]\or\faStar\faStar\faStar\fi}}
\newcounter{exer}
\newcommand{\ex}[1]{\par\needspace{4\baselineskip}\medskip\refstepcounter{exer}\noindent{\color{navy}\bfseries\theexer.}\ \ {\small\estrelles{#1}}\ \ }

\newcommand{\seccio}[2]{\par\needspace{16\baselineskip}\bigskip
  \begin{tcolorbox}[colback=navy,colframe=navy,coltext=white,arc=2mm,boxrule=0pt,left=3mm,right=3mm,top=1.2mm,bottom=1.2mm]
  {\large\bfseries #1\quad #2}
  \end{tcolorbox}\vspace{-2pt}}

\newcommand{\exf}[1]{\par\medskip\refstepcounter{exer}\noindent{\color{navy}\bfseries\theexer.}\ \ {\small\estrelles{#1}}\ \ }
\newcolumntype{C}[1]{>{\centering\arraybackslash}p{#1}}
% quadre de recordatori (només a la fitxa d'exercicis)
\newcommand{\recorda}[1]{\par\smallskip\begin{tcolorbox}[colback=accent!7,colframe=accent!55,arc=2mm,boxrule=0.5pt,left=2.5mm,right=2.5mm,top=1mm,bottom=1mm]\small{\color{accent}\faLightbulb[regular]}\ \textbf{Recorda:}\ #1\end{tcolorbox}}
\usetikzlibrary{calc}
\usepackage{enumitem}
\definecolor{violet}{HTML}{7C3AED}
\definecolor{orange}{HTML}{EA8A1F}
\newcommand{\res}[1]{#1}\newcommand{\rest}[1]{#1}

\def\fmneed{11}% que el títol d'una secció no quede sol al final del full
% =====================================================================
% Capçalera FentMates de les fitxes: un color per a cada curs.
% Els colors són els de la pàgina inicial (LEVELS a assets/js/materials.js);
% el logo ve de _eines/logo/pdf/ (python3 _eines/logo/genera.py).
%
% Ús, al final del preàmbul de la unitat:
%   % =====================================================================
% Capçalera FentMates de les fitxes: un color per a cada curs.
% Els colors són els de la pàgina inicial (LEVELS a assets/js/materials.js);
% el logo ve de _eines/logo/pdf/ (python3 _eines/logo/genera.py).
%
% Ús, al final del preàmbul de la unitat:
%   % =====================================================================
% Capçalera FentMates de les fitxes: un color per a cada curs.
% Els colors són els de la pàgina inicial (LEVELS a assets/js/materials.js);
% el logo ve de _eines/logo/pdf/ (python3 _eines/logo/genera.py).
%
% Ús, al final del preàmbul de la unitat:
%   \input{../comu/fentmates.tex}
%   \fmcurs{1eso}
% i, al començament del document:
%   \capcalera{Matemàtiques}{1r ESO}{Fitxa d'exercicis · Unitat 2}{Divisibilitat}
% Opcional abans de \input: \def\fmneed{14} (espai mínim abans de cada secció).
% =====================================================================
\makeatletter
\def\fm@logodir{../../logo/pdf}
\providecommand{\fmneed}{10}
\def\fm@cursnom{}

% clau del curs, color i color del text damunt del color
\newcommand{\fm@curs}[3]{\@namedef{fm@color@#1}{#2}\@namedef{fm@text@#1}{#3}}
\fm@curs{1eso}{E6B84A}{navy}
\fm@curs{2eso}{3B82F6}{white}
\fm@curs{3eso}{FF7A50}{white}
\fm@curs{4eso-a}{A78BFA}{white}
\fm@curs{4eso-b}{6D28D9}{white}
\fm@curs{1batx-mi}{9B1C1C}{white}
\fm@curs{1batx-ccss}{EF6B6B}{white}
\fm@curs{2batx-mii}{1E6B47}{white}
\fm@curs{2batx-ccss}{2BB3A3}{white}

\newcommand{\fmcurs}[1]{%
  \@ifundefined{fm@color@#1}{\PackageError{fentmates}{Curs desconegut: #1}{}}{}%
  \edef\fm@tmp{\noexpand\definecolor{curs}{HTML}{\@nameuse{fm@color@#1}}%
               \noexpand\colorlet{curstext}{\@nameuse{fm@text@#1}}}\fm@tmp
  \colorlet{cursfosc}{curs!72!black}%
  \def\fm@logo{\fm@logodir/vertical-#1.pdf}%
  \def\fm@marca{\fm@logodir/marca-#1.pdf}}

% ---------- capçalera de la primera pàgina ----------
% #1 matèria (petit), #2 curs, #3 tipus de document i unitat, #4 títol
\newcommand{\capcalera}[4]{%
  \gdef\fm@cursnom{#2}%
  \par\noindent
  \begin{tcolorbox}[colback=curs,colframe=curs,boxrule=0pt,arc=2.2mm,width=3.5cm,height=2.15cm,
      valign=center,left=3mm,right=2mm,top=1mm,bottom=1mm,nobeforeafter,box align=center]
    \color{curstext}{\footnotesize #1}\par\vspace{1pt}{\fontsize{19}{21}\selectfont\bfseries #2\par}%
  \end{tcolorbox}\hspace{5mm}%
  \begin{minipage}[c]{\dimexpr\linewidth-3.5cm-5mm-2.9cm\relax}
    {\color{cursfosc}\small\bfseries\MakeUppercase{#3}}\par\vspace{2pt}
    {\color{navy}\fontsize{20}{23}\selectfont\bfseries #4\par}
  \end{minipage}\hfill
  \begin{minipage}[c]{2.6cm}\raggedleft\includegraphics[width=2.6cm]{\fm@logo}\end{minipage}\par
  \vspace{2pt}}

% ---------- barres de secció del color del curs ----------
\renewcommand{\seccio}[2]{\par\needspace{\fmneed\baselineskip}\bigskip
  \begin{tcolorbox}[colback=curs,colframe=curs,coltext=curstext,arc=2mm,boxrule=0pt,left=3mm,right=3mm,top=1.2mm,bottom=1.2mm]
  {\large\bfseries #1\quad #2}
  \end{tcolorbox}\vspace{-2pt}}

% ---------- estrelles de dificultat del color del curs (a Batxillerat no se'n posen) ----------
\@ifundefined{estrelles}{}{%
  \renewcommand{\estrelles}[1]{{\color{curs}\ifcase#1\or\faStar\faStar[regular]\faStar[regular]\or\faStar\faStar\faStar[regular]\or\faStar\faStar\faStar\fi}}}

% ---------- capçalera i peu de les pàgines ----------
% primera pàgina: logos institucionals al peu
\fancypagestyle{primera}{\fancyhf{}\renewcommand{\headrulewidth}{0pt}
  \fancyfoot[L]{\raisebox{-0.32\height}{\includegraphics[height=0.78cm]{logo-gva.png}}\hspace{5mm}%
    \raisebox{-0.32\height}{\includegraphics[height=0.78cm]{logo-vermellar.png}}}
  \fancyfoot[R]{\scriptsize\color{muted}\href{https://creativecommons.org/licenses/by-nc-sa/4.0/deed.ca}{\ccbyncsa}\ \ CC BY-NC-SA 4.0 · fentmates.org\qquad\small\thepage}}
% la resta: el símbol a l'esquerra i el curs a la dreta, amb el títol al centre com abans
\fancyhead[L]{\raisebox{-0.3\height}{\includegraphics[height=0.62cm]{\fm@marca}}}
\fancyhead[R]{\small\bfseries\color{cursfosc}\fm@cursnom}
\renewcommand{\headrule}{\color{curs!45}\hrule height \headrulewidth}
\makeatother

%   \fmcurs{1eso}
% i, al començament del document:
%   \capcalera{Matemàtiques}{1r ESO}{Fitxa d'exercicis · Unitat 2}{Divisibilitat}
% Opcional abans de \input: \def\fmneed{14} (espai mínim abans de cada secció).
% =====================================================================
\makeatletter
\def\fm@logodir{../../logo/pdf}
\providecommand{\fmneed}{10}
\def\fm@cursnom{}

% clau del curs, color i color del text damunt del color
\newcommand{\fm@curs}[3]{\@namedef{fm@color@#1}{#2}\@namedef{fm@text@#1}{#3}}
\fm@curs{1eso}{E6B84A}{navy}
\fm@curs{2eso}{3B82F6}{white}
\fm@curs{3eso}{FF7A50}{white}
\fm@curs{4eso-a}{A78BFA}{white}
\fm@curs{4eso-b}{6D28D9}{white}
\fm@curs{1batx-mi}{9B1C1C}{white}
\fm@curs{1batx-ccss}{EF6B6B}{white}
\fm@curs{2batx-mii}{1E6B47}{white}
\fm@curs{2batx-ccss}{2BB3A3}{white}

\newcommand{\fmcurs}[1]{%
  \@ifundefined{fm@color@#1}{\PackageError{fentmates}{Curs desconegut: #1}{}}{}%
  \edef\fm@tmp{\noexpand\definecolor{curs}{HTML}{\@nameuse{fm@color@#1}}%
               \noexpand\colorlet{curstext}{\@nameuse{fm@text@#1}}}\fm@tmp
  \colorlet{cursfosc}{curs!72!black}%
  \def\fm@logo{\fm@logodir/vertical-#1.pdf}%
  \def\fm@marca{\fm@logodir/marca-#1.pdf}}

% ---------- capçalera de la primera pàgina ----------
% #1 matèria (petit), #2 curs, #3 tipus de document i unitat, #4 títol
\newcommand{\capcalera}[4]{%
  \gdef\fm@cursnom{#2}%
  \par\noindent
  \begin{tcolorbox}[colback=curs,colframe=curs,boxrule=0pt,arc=2.2mm,width=3.5cm,height=2.15cm,
      valign=center,left=3mm,right=2mm,top=1mm,bottom=1mm,nobeforeafter,box align=center]
    \color{curstext}{\footnotesize #1}\par\vspace{1pt}{\fontsize{19}{21}\selectfont\bfseries #2\par}%
  \end{tcolorbox}\hspace{5mm}%
  \begin{minipage}[c]{\dimexpr\linewidth-3.5cm-5mm-2.9cm\relax}
    {\color{cursfosc}\small\bfseries\MakeUppercase{#3}}\par\vspace{2pt}
    {\color{navy}\fontsize{20}{23}\selectfont\bfseries #4\par}
  \end{minipage}\hfill
  \begin{minipage}[c]{2.6cm}\raggedleft\includegraphics[width=2.6cm]{\fm@logo}\end{minipage}\par
  \vspace{2pt}}

% ---------- barres de secció del color del curs ----------
\renewcommand{\seccio}[2]{\par\needspace{\fmneed\baselineskip}\bigskip
  \begin{tcolorbox}[colback=curs,colframe=curs,coltext=curstext,arc=2mm,boxrule=0pt,left=3mm,right=3mm,top=1.2mm,bottom=1.2mm]
  {\large\bfseries #1\quad #2}
  \end{tcolorbox}\vspace{-2pt}}

% ---------- estrelles de dificultat del color del curs (a Batxillerat no se'n posen) ----------
\@ifundefined{estrelles}{}{%
  \renewcommand{\estrelles}[1]{{\color{curs}\ifcase#1\or\faStar\faStar[regular]\faStar[regular]\or\faStar\faStar\faStar[regular]\or\faStar\faStar\faStar\fi}}}

% ---------- capçalera i peu de les pàgines ----------
% primera pàgina: logos institucionals al peu
\fancypagestyle{primera}{\fancyhf{}\renewcommand{\headrulewidth}{0pt}
  \fancyfoot[L]{\raisebox{-0.32\height}{\includegraphics[height=0.78cm]{logo-gva.png}}\hspace{5mm}%
    \raisebox{-0.32\height}{\includegraphics[height=0.78cm]{logo-vermellar.png}}}
  \fancyfoot[R]{\scriptsize\color{muted}\href{https://creativecommons.org/licenses/by-nc-sa/4.0/deed.ca}{\ccbyncsa}\ \ CC BY-NC-SA 4.0 · fentmates.org\qquad\small\thepage}}
% la resta: el símbol a l'esquerra i el curs a la dreta, amb el títol al centre com abans
\fancyhead[L]{\raisebox{-0.3\height}{\includegraphics[height=0.62cm]{\fm@marca}}}
\fancyhead[R]{\small\bfseries\color{cursfosc}\fm@cursnom}
\renewcommand{\headrule}{\color{curs!45}\hrule height \headrulewidth}
\makeatother

%   \fmcurs{1eso}
% i, al començament del document:
%   \capcalera{Matemàtiques}{1r ESO}{Fitxa d'exercicis · Unitat 2}{Divisibilitat}
% Opcional abans de \input: \def\fmneed{14} (espai mínim abans de cada secció).
% =====================================================================
\makeatletter
\def\fm@logodir{../../logo/pdf}
\providecommand{\fmneed}{10}
\def\fm@cursnom{}

% clau del curs, color i color del text damunt del color
\newcommand{\fm@curs}[3]{\@namedef{fm@color@#1}{#2}\@namedef{fm@text@#1}{#3}}
\fm@curs{1eso}{E6B84A}{navy}
\fm@curs{2eso}{3B82F6}{white}
\fm@curs{3eso}{FF7A50}{white}
\fm@curs{4eso-a}{A78BFA}{white}
\fm@curs{4eso-b}{6D28D9}{white}
\fm@curs{1batx-mi}{9B1C1C}{white}
\fm@curs{1batx-ccss}{EF6B6B}{white}
\fm@curs{2batx-mii}{1E6B47}{white}
\fm@curs{2batx-ccss}{2BB3A3}{white}

\newcommand{\fmcurs}[1]{%
  \@ifundefined{fm@color@#1}{\PackageError{fentmates}{Curs desconegut: #1}{}}{}%
  \edef\fm@tmp{\noexpand\definecolor{curs}{HTML}{\@nameuse{fm@color@#1}}%
               \noexpand\colorlet{curstext}{\@nameuse{fm@text@#1}}}\fm@tmp
  \colorlet{cursfosc}{curs!72!black}%
  \def\fm@logo{\fm@logodir/vertical-#1.pdf}%
  \def\fm@marca{\fm@logodir/marca-#1.pdf}}

% ---------- capçalera de la primera pàgina ----------
% #1 matèria (petit), #2 curs, #3 tipus de document i unitat, #4 títol
\newcommand{\capcalera}[4]{%
  \gdef\fm@cursnom{#2}%
  \par\noindent
  \begin{tcolorbox}[colback=curs,colframe=curs,boxrule=0pt,arc=2.2mm,width=3.5cm,height=2.15cm,
      valign=center,left=3mm,right=2mm,top=1mm,bottom=1mm,nobeforeafter,box align=center]
    \color{curstext}{\footnotesize #1}\par\vspace{1pt}{\fontsize{19}{21}\selectfont\bfseries #2\par}%
  \end{tcolorbox}\hspace{5mm}%
  \begin{minipage}[c]{\dimexpr\linewidth-3.5cm-5mm-2.9cm\relax}
    {\color{cursfosc}\small\bfseries\MakeUppercase{#3}}\par\vspace{2pt}
    {\color{navy}\fontsize{20}{23}\selectfont\bfseries #4\par}
  \end{minipage}\hfill
  \begin{minipage}[c]{2.6cm}\raggedleft\includegraphics[width=2.6cm]{\fm@logo}\end{minipage}\par
  \vspace{2pt}}

% ---------- barres de secció del color del curs ----------
\renewcommand{\seccio}[2]{\par\needspace{\fmneed\baselineskip}\bigskip
  \begin{tcolorbox}[colback=curs,colframe=curs,coltext=curstext,arc=2mm,boxrule=0pt,left=3mm,right=3mm,top=1.2mm,bottom=1.2mm]
  {\large\bfseries #1\quad #2}
  \end{tcolorbox}\vspace{-2pt}}

% ---------- estrelles de dificultat del color del curs (a Batxillerat no se'n posen) ----------
\@ifundefined{estrelles}{}{%
  \renewcommand{\estrelles}[1]{{\color{curs}\ifcase#1\or\faStar\faStar[regular]\faStar[regular]\or\faStar\faStar\faStar[regular]\or\faStar\faStar\faStar\fi}}}

% ---------- capçalera i peu de les pàgines ----------
% primera pàgina: logos institucionals al peu
\fancypagestyle{primera}{\fancyhf{}\renewcommand{\headrulewidth}{0pt}
  \fancyfoot[L]{\raisebox{-0.32\height}{\includegraphics[height=0.78cm]{logo-gva.png}}\hspace{5mm}%
    \raisebox{-0.32\height}{\includegraphics[height=0.78cm]{logo-vermellar.png}}}
  \fancyfoot[R]{\scriptsize\color{muted}\href{https://creativecommons.org/licenses/by-nc-sa/4.0/deed.ca}{\ccbyncsa}\ \ CC BY-NC-SA 4.0 · fentmates.org\qquad\small\thepage}}
% la resta: el símbol a l'esquerra i el curs a la dreta, amb el títol al centre com abans
\fancyhead[L]{\raisebox{-0.3\height}{\includegraphics[height=0.62cm]{\fm@marca}}}
\fancyhead[R]{\small\bfseries\color{cursfosc}\fm@cursnom}
\renewcommand{\headrule}{\color{curs!45}\hrule height \headrulewidth}
\makeatother

\fmcurs{4eso-a}

\begin{document}
\thispagestyle{primera}

\capcalera{Matemàtiques A}{4t ESO}{Fitxa d'exercicis · Unitat 7}{\raggedright Funcions elementals}
\vspace{10pt}
Nom i cognoms: \rule{8.4cm}{0.4pt}\hfill Grup: \rule{1.6cm}{0.4pt}\hfill Data: \rule{2.2cm}{0.4pt}
\vspace{8pt}

\begin{tcolorbox}[colback=soft,colframe=soft,arc=2mm,left=3mm,right=3mm,top=1.5mm,bottom=1.5mm]
\small Exercicis graduats per nivells de dificultat:\quad \estrelles{1}\ bàsic \qquad \estrelles{2}\ mitjà \qquad \estrelles{3}\ avançat.\\[2pt] Resol-los al quadern mostrant tots els passos i dibuixa les gràfiques en paper quadriculat. Escriu els intervals com als apunts: $(a,b)$ obert i $[a,b]$ tancat. Les seccions segueixen l'ordre dels apunts de la unitat (7.1--7.9).
\end{tcolorbox}

% =====================================================
\seccio{1}{Funcions lineals}

\exf{1} Dibuixa les rectes següents i digues de quin tipus és cadascuna (de proporcionalitat directa, afí o constant):\begin{tasks}(3)
  \task $f(x)=2x+1$
  \task $f(x)=-x-2$
  \task $f(x)=3x$
  \task $f(x)=-2x$
  \task $f(x)=-2$
  \task $f(x)=\frac12x+2$
\end{tasks}\recorda{per a dibuixar una recta n'hi ha prou amb dos punts: fes una taula de valors amb dos valors de $x$.}

\ex{1} Sense dibuixar-les, indica el pendent, l'ordenada en l'origen i el tipus de cada recta, i digues si és creixent, decreixent o constant:\begin{tasks}(4)
  \task $y=3x-5$
  \task $y=-\dfrac x2+4$
  \task $y=7$
  \task $y=-4x$
\end{tasks}

\ex{1} Escriu la funció que relaciona les dues magnituds i digues de quin tipus és:\begin{tasks}(1)
  \task El preu $y$ de $x$ kg de taronges a 1,20 € el quilo.
  \task El que paguem cada mes, $y$, en una tarifa plana de 25 € al mes, segons les hores $x$ que ens connectem.
  \task El preu $y$ d'una carrera de taxi de $x$ km si cobra 3 € per la baixada de bandera i 1 € per km.
  \task L'altura $y$ (en cm) d'una espelma de 20 cm que es consumeix 2 cm cada hora, al cap de $x$ hores.
\end{tasks}

\ex{2} Digues si cada taula correspon a una funció de proporcionalitat directa o afí i escriu-ne l'expressió:\begin{tasks}(2)
  \task \begin{center}\renewcommand{\arraystretch}{1.25}\begin{tabular}{|c|C{1cm}|C{1cm}|C{1cm}|C{1cm}|}\hline $x$ & $1$ & $2$ & $3$ & $5$\\\hline $y$ & $3$ & $6$ & $9$ & $15$\\\hline\end{tabular}\end{center}
  \task \begin{center}\renewcommand{\arraystretch}{1.25}\begin{tabular}{|c|C{1cm}|C{1cm}|C{1cm}|C{1cm}|}\hline $x$ & $0$ & $2$ & $4$ & $6$\\\hline $y$ & $3$ & $7$ & $11$ & $15$\\\hline\end{tabular}\end{center}
\end{tasks}

\ex{2} Representa en uns mateixos eixos $y=2x$, $y=2x+3$ i $y=2x-1$. Què tenen en comú? En què es diferencien?


% =====================================================
\seccio{2}{Equació punt-pendent}

\exf{1} Calcula el pendent de la recta que passa pels punts:\begin{tasks}(4)
  \task $(1,\,2)$ i $(3,\,8)$
  \task $(-2,\,5)$ i $(4,\,-1)$
  \task $(0,\,4)$ i $(6,\,1)$
  \task $(-3,\,-2)$ i $(1,\,6)$
\end{tasks}\recorda{$m=\dfrac{y_2-y_1}{x_2-x_1}$.}

\ex{1} Escriu l'equació punt-pendent i l'equació explícita ($y=mx+n$) de la recta:\begin{tasks}(2)
  \task De pendent $3$ que passa per $(2,\,1)$.
  \task De pendent $-2$ que passa per $(-1,\,4)$.
  \task De pendent $\frac12$ que passa per $(4,\,3)$.
  \task De pendent $0$ que passa per $(5,\,-2)$.
\end{tasks}\recorda{$y=y_1+m(x-x_1)$.}

\ex{1} Troba l'equació de la recta que passa pels punts:\begin{tasks}(4)
  \task $(1,\,2)$ i $(3,\,8)$
  \task $(-2,\,5)$ i $(4,\,-1)$
  \task $(0,\,-3)$ i $(2,\,1)$
  \task $(-1,\,-4)$ i $(3,\,4)$
\end{tasks}

\ex{2} Troba l'expressió de cada recta a partir de la gràfica i digues de quin tipus és:

\begin{center}\begin{tikzpicture}[x=0.636cm,y=0.636cm,baseline={(current bounding box.center)},font=\scriptsize]
\draw[muted!25,line width=.3pt] (-5,-4)--(-5,4);
\draw[muted!25,line width=.3pt] (-4,-4)--(-4,4);
\draw[muted!25,line width=.3pt] (-3,-4)--(-3,4);
\draw[muted!25,line width=.3pt] (-2,-4)--(-2,4);
\draw[muted!25,line width=.3pt] (-1,-4)--(-1,4);
\draw[muted!25,line width=.3pt] (0,-4)--(0,4);
\draw[muted!25,line width=.3pt] (1,-4)--(1,4);
\draw[muted!25,line width=.3pt] (2,-4)--(2,4);
\draw[muted!25,line width=.3pt] (3,-4)--(3,4);
\draw[muted!25,line width=.3pt] (4,-4)--(4,4);
\draw[muted!25,line width=.3pt] (5,-4)--(5,4);
\draw[muted!25,line width=.3pt] (6,-4)--(6,4);
\draw[muted!25,line width=.3pt] (-5,-4)--(6,-4);
\draw[muted!25,line width=.3pt] (-5,-3)--(6,-3);
\draw[muted!25,line width=.3pt] (-5,-2)--(6,-2);
\draw[muted!25,line width=.3pt] (-5,-1)--(6,-1);
\draw[muted!25,line width=.3pt] (-5,0)--(6,0);
\draw[muted!25,line width=.3pt] (-5,1)--(6,1);
\draw[muted!25,line width=.3pt] (-5,2)--(6,2);
\draw[muted!25,line width=.3pt] (-5,3)--(6,3);
\draw[muted!25,line width=.3pt] (-5,4)--(6,4);
\draw[-latex,muted] (-5,0)--(6,0);
\draw[-latex,muted] (0,-4)--(0,4);
\node[below=-1pt,text=muted] at (-5,0) {$-5$};
\node[below=-1pt,text=muted] at (-4,0) {$-4$};
\node[below=-1pt,text=muted] at (-3,0) {$-3$};
\node[below=-1pt,text=muted] at (-2,0) {$-2$};
\node[below=-1pt,text=muted] at (-1,0) {$-1$};
\node[below=-1pt,text=muted] at (1,0) {$1$};
\node[below=-1pt,text=muted] at (2,0) {$2$};
\node[below=-1pt,text=muted] at (3,0) {$3$};
\node[below=-1pt,text=muted] at (4,0) {$4$};
\node[below=-1pt,text=muted] at (5,0) {$5$};
\node[below=-1pt,text=muted] at (6,0) {$6$};
\node[left=-1pt,text=muted] at (0,-4) {$-4$};
\node[left=-1pt,text=muted] at (0,-3) {$-3$};
\node[left=-1pt,text=muted] at (0,-2) {$-2$};
\node[left=-1pt,text=muted] at (0,-1) {$-1$};
\node[left=-1pt,text=muted] at (0,1) {$1$};
\node[left=-1pt,text=muted] at (0,2) {$2$};
\node[left=-1pt,text=muted] at (0,3) {$3$};
\node[left=-1pt,text=muted] at (0,4) {$4$};
\begin{scope}\clip (-5,-4) rectangle (6,4);
\draw[green!55!black,line width=1.1pt,line join=round] (-5,2)--(-4.912,2)--(-4.824,2)--(-4.736,2)--(-4.648,2)--(-4.56,2)--(-4.472,2)--(-4.384,2)--(-4.296,2)--(-4.208,2)--(-4.12,2)--(-4.032,2)--(-3.944,2)--(-3.856,2)--(-3.768,2)--(-3.68,2)--(-3.592,2)--(-3.504,2)--(-3.416,2)--(-3.328,2)--(-3.24,2)--(-3.152,2)--(-3.064,2)--(-2.976,2)--(-2.888,2)--(-2.8,2)--(-2.712,2)--(-2.624,2)--(-2.536,2)--(-2.448,2)--(-2.36,2)--(-2.272,2)--(-2.184,2)--(-2.096,2)--(-2.008,2)--(-1.92,2)--(-1.832,2)--(-1.744,2)--(-1.656,2)--(-1.568,2)--(-1.48,2)--(-1.392,2)--(-1.304,2)--(-1.216,2)--(-1.128,2)--(-1.04,2)--(-0.952,2)--(-0.864,2)--(-0.776,2)--(-0.688,2)--(-0.6,2)--(-0.512,2)--(-0.424,2)--(-0.336,2)--(-0.248,2)--(-0.16,2)--(-0.072,2)--(0.016,2)--(0.104,2)--(0.192,2)--(0.28,2)--(0.368,2)--(0.456,2)--(0.544,2)--(0.632,2)--(0.72,2)--(0.808,2)--(0.896,2)--(0.984,2)--(1.072,2)--(1.16,2)--(1.248,2)--(1.336,2)--(1.424,2)--(1.512,2)--(1.6,2)--(1.688,2)--(1.776,2)--(1.864,2)--(1.952,2)--(2.04,2)--(2.128,2)--(2.216,2)--(2.304,2)--(2.392,2)--(2.48,2)--(2.568,2)--(2.656,2)--(2.744,2)--(2.832,2)--(2.92,2)--(3.008,2)--(3.096,2)--(3.184,2)--(3.272,2)--(3.36,2)--(3.448,2)--(3.536,2)--(3.624,2)--(3.712,2)--(3.8,2)--(3.888,2)--(3.976,2)--(4.064,2)--(4.152,2)--(4.24,2)--(4.328,2)--(4.416,2)--(4.504,2)--(4.592,2)--(4.68,2)--(4.768,2)--(4.856,2)--(4.944,2)--(5.032,2)--(5.12,2)--(5.208,2)--(5.296,2)--(5.384,2)--(5.472,2)--(5.56,2)--(5.648,2)--(5.736,2)--(5.824,2)--(5.912,2)--(6,2);
\draw[blue!70!black,line width=1.1pt,line join=round] (-5,-9)--(-4.912,-8.824)--(-4.824,-8.648)--(-4.736,-8.472)--(-4.648,-8.296)--(-4.56,-8.12)--(-4.472,-7.944)--(-4.384,-7.768)--(-4.296,-7.592)--(-4.208,-7.416)--(-4.12,-7.24)--(-4.032,-7.064)--(-3.944,-6.888)--(-3.856,-6.712)--(-3.768,-6.536)--(-3.68,-6.36)--(-3.592,-6.184)--(-3.504,-6.008)--(-3.416,-5.832)--(-3.328,-5.656)--(-3.24,-5.48)--(-3.152,-5.304)--(-3.064,-5.128)--(-2.976,-4.952)--(-2.888,-4.776)--(-2.8,-4.6)--(-2.712,-4.424)--(-2.624,-4.248)--(-2.536,-4.072)--(-2.448,-3.896)--(-2.36,-3.72)--(-2.272,-3.544)--(-2.184,-3.368)--(-2.096,-3.192)--(-2.008,-3.016)--(-1.92,-2.84)--(-1.832,-2.664)--(-1.744,-2.488)--(-1.656,-2.312)--(-1.568,-2.136)--(-1.48,-1.96)--(-1.392,-1.784)--(-1.304,-1.608)--(-1.216,-1.432)--(-1.128,-1.256)--(-1.04,-1.08)--(-0.952,-0.904)--(-0.864,-0.728)--(-0.776,-0.552)--(-0.688,-0.376)--(-0.6,-0.2)--(-0.512,-0.024)--(-0.424,0.152)--(-0.336,0.328)--(-0.248,0.504)--(-0.16,0.68)--(-0.072,0.856)--(0.016,1.032)--(0.104,1.208)--(0.192,1.384)--(0.28,1.56)--(0.368,1.736)--(0.456,1.912)--(0.544,2.088)--(0.632,2.264)--(0.72,2.44)--(0.808,2.616)--(0.896,2.792)--(0.984,2.968)--(1.072,3.144)--(1.16,3.32)--(1.248,3.496)--(1.336,3.672)--(1.424,3.848)--(1.512,4.024)--(1.6,4.2)--(1.688,4.376)--(1.776,4.552)--(1.864,4.728)--(1.952,4.904)--(2.04,5.08)--(2.128,5.256)--(2.216,5.432)--(2.304,5.608)--(2.392,5.784)--(2.48,5.96)--(2.568,6.136)--(2.656,6.312)--(2.744,6.488)--(2.832,6.664)--(2.92,6.84)--(3.008,7.016)--(3.096,7.192)--(3.184,7.368)--(3.272,7.544)--(3.36,7.72)--(3.448,7.896)--(3.536,8.072)--(3.624,8.248)--(3.712,8.424)--(3.8,8.6)--(3.888,8.776)--(3.976,8.952)--(4.064,9.128)--(4.152,9.304)--(4.24,9.48)--(4.328,9.656)--(4.416,9.832)--(4.504,10.008)--(4.592,10.184)--(4.68,10.36)--(4.768,10.536)--(4.856,10.712)--(4.944,10.888)--(5.032,11.064)--(5.12,11.24)--(5.208,11.416)--(5.296,11.592)--(5.384,11.768)--(5.472,11.944)--(5.56,12)--(5.648,12)--(5.736,12)--(5.824,12)--(5.912,12)--(6,12);
\draw[accent,line width=1.1pt,line join=round] (-5,4)--(-4.912,3.912)--(-4.824,3.824)--(-4.736,3.736)--(-4.648,3.648)--(-4.56,3.56)--(-4.472,3.472)--(-4.384,3.384)--(-4.296,3.296)--(-4.208,3.208)--(-4.12,3.12)--(-4.032,3.032)--(-3.944,2.944)--(-3.856,2.856)--(-3.768,2.768)--(-3.68,2.68)--(-3.592,2.592)--(-3.504,2.504)--(-3.416,2.416)--(-3.328,2.328)--(-3.24,2.24)--(-3.152,2.152)--(-3.064,2.064)--(-2.976,1.976)--(-2.888,1.888)--(-2.8,1.8)--(-2.712,1.712)--(-2.624,1.624)--(-2.536,1.536)--(-2.448,1.448)--(-2.36,1.36)--(-2.272,1.272)--(-2.184,1.184)--(-2.096,1.096)--(-2.008,1.008)--(-1.92,0.92)--(-1.832,0.832)--(-1.744,0.744)--(-1.656,0.656)--(-1.568,0.568)--(-1.48,0.48)--(-1.392,0.392)--(-1.304,0.304)--(-1.216,0.216)--(-1.128,0.128)--(-1.04,0.04)--(-0.952,-0.048)--(-0.864,-0.136)--(-0.776,-0.224)--(-0.688,-0.312)--(-0.6,-0.4)--(-0.512,-0.488)--(-0.424,-0.576)--(-0.336,-0.664)--(-0.248,-0.752)--(-0.16,-0.84)--(-0.072,-0.928)--(0.016,-1.016)--(0.104,-1.104)--(0.192,-1.192)--(0.28,-1.28)--(0.368,-1.368)--(0.456,-1.456)--(0.544,-1.544)--(0.632,-1.632)--(0.72,-1.72)--(0.808,-1.808)--(0.896,-1.896)--(0.984,-1.984)--(1.072,-2.072)--(1.16,-2.16)--(1.248,-2.248)--(1.336,-2.336)--(1.424,-2.424)--(1.512,-2.512)--(1.6,-2.6)--(1.688,-2.688)--(1.776,-2.776)--(1.864,-2.864)--(1.952,-2.952)--(2.04,-3.04)--(2.128,-3.128)--(2.216,-3.216)--(2.304,-3.304)--(2.392,-3.392)--(2.48,-3.48)--(2.568,-3.568)--(2.656,-3.656)--(2.744,-3.744)--(2.832,-3.832)--(2.92,-3.92)--(3.008,-4.008)--(3.096,-4.096)--(3.184,-4.184)--(3.272,-4.272)--(3.36,-4.36)--(3.448,-4.448)--(3.536,-4.536)--(3.624,-4.624)--(3.712,-4.712)--(3.8,-4.8)--(3.888,-4.888)--(3.976,-4.976)--(4.064,-5.064)--(4.152,-5.152)--(4.24,-5.24)--(4.328,-5.328)--(4.416,-5.416)--(4.504,-5.504)--(4.592,-5.592)--(4.68,-5.68)--(4.768,-5.768)--(4.856,-5.856)--(4.944,-5.944)--(5.032,-6.032)--(5.12,-6.12)--(5.208,-6.208)--(5.296,-6.296)--(5.384,-6.384)--(5.472,-6.472)--(5.56,-6.56)--(5.648,-6.648)--(5.736,-6.736)--(5.824,-6.824)--(5.912,-6.912)--(6,-7);
\draw[orange!85!black,line width=1.1pt,line join=round] (-5,-2.5)--(-4.912,-2.456)--(-4.824,-2.412)--(-4.736,-2.368)--(-4.648,-2.324)--(-4.56,-2.28)--(-4.472,-2.236)--(-4.384,-2.192)--(-4.296,-2.148)--(-4.208,-2.104)--(-4.12,-2.06)--(-4.032,-2.016)--(-3.944,-1.972)--(-3.856,-1.928)--(-3.768,-1.884)--(-3.68,-1.84)--(-3.592,-1.796)--(-3.504,-1.752)--(-3.416,-1.708)--(-3.328,-1.664)--(-3.24,-1.62)--(-3.152,-1.576)--(-3.064,-1.532)--(-2.976,-1.488)--(-2.888,-1.444)--(-2.8,-1.4)--(-2.712,-1.356)--(-2.624,-1.312)--(-2.536,-1.268)--(-2.448,-1.224)--(-2.36,-1.18)--(-2.272,-1.136)--(-2.184,-1.092)--(-2.096,-1.048)--(-2.008,-1.004)--(-1.92,-0.96)--(-1.832,-0.916)--(-1.744,-0.872)--(-1.656,-0.828)--(-1.568,-0.784)--(-1.48,-0.74)--(-1.392,-0.696)--(-1.304,-0.652)--(-1.216,-0.608)--(-1.128,-0.564)--(-1.04,-0.52)--(-0.952,-0.476)--(-0.864,-0.432)--(-0.776,-0.388)--(-0.688,-0.344)--(-0.6,-0.3)--(-0.512,-0.256)--(-0.424,-0.212)--(-0.336,-0.168)--(-0.248,-0.124)--(-0.16,-0.08)--(-0.072,-0.036)--(0.016,0.008)--(0.104,0.052)--(0.192,0.096)--(0.28,0.14)--(0.368,0.184)--(0.456,0.228)--(0.544,0.272)--(0.632,0.316)--(0.72,0.36)--(0.808,0.404)--(0.896,0.448)--(0.984,0.492)--(1.072,0.536)--(1.16,0.58)--(1.248,0.624)--(1.336,0.668)--(1.424,0.712)--(1.512,0.756)--(1.6,0.8)--(1.688,0.844)--(1.776,0.888)--(1.864,0.932)--(1.952,0.976)--(2.04,1.02)--(2.128,1.064)--(2.216,1.108)--(2.304,1.152)--(2.392,1.196)--(2.48,1.24)--(2.568,1.284)--(2.656,1.328)--(2.744,1.372)--(2.832,1.416)--(2.92,1.46)--(3.008,1.504)--(3.096,1.548)--(3.184,1.592)--(3.272,1.636)--(3.36,1.68)--(3.448,1.724)--(3.536,1.768)--(3.624,1.812)--(3.712,1.856)--(3.8,1.9)--(3.888,1.944)--(3.976,1.988)--(4.064,2.032)--(4.152,2.076)--(4.24,2.12)--(4.328,2.164)--(4.416,2.208)--(4.504,2.252)--(4.592,2.296)--(4.68,2.34)--(4.768,2.384)--(4.856,2.428)--(4.944,2.472)--(5.032,2.516)--(5.12,2.56)--(5.208,2.604)--(5.296,2.648)--(5.384,2.692)--(5.472,2.736)--(5.56,2.78)--(5.648,2.824)--(5.736,2.868)--(5.824,2.912)--(5.912,2.956)--(6,3);
\node[text=green!55!black] at (5.6,2.45) {\normalsize$\boldsymbol{f}$};
\node[text=blue!70!black] at (1.9,3.5) {\normalsize$\boldsymbol{g}$};
\node[text=accent] at (-4.5,3) {\normalsize$\boldsymbol{h}$};
\node[text=orange!85!black] at (5.5,3.3) {\normalsize$\boldsymbol{i}$};
\end{scope}\end{tikzpicture}\end{center}


\ex{2} Troba l'equació de la recta:\begin{tasks}(1)
  \task Paral·lela a $y=2x-7$ que passa per $(0,\,4)$.
  \task Paral·lela a $y=-3x+1$ que passa per $(2,\,-1)$.
  \task Paral·lela a l'eix $OX$ que passa per $(3,\,5)$.
\end{tasks}

\ex{2} Comprova si els tres punts estan alineats (si són de la mateixa recta):\begin{tasks}(2)
  \task $(1,\,1)$, $(3,\,5)$ i $(5,\,9)$
  \task $(0,\,2)$, $(2,\,5)$ i $(4,\,9)$
\end{tasks}

\ex{2} Representa en uns mateixos eixos $f(x)=2x-1$ i $g(x)=-x+5$ i troba'n el punt de tall, gràficament i resolent una equació.

\ex{3} Una recta passa pels punts $A(2,\,3)$ i $B(6,\,5)$. Troba'n l'equació, els punts de tall amb els eixos i comprova si el punt $(10,\,7)$ és de la recta.


% =====================================================
\seccio{3}{Funció quadràtica}

\exf{1} Sense dibuixar-les, indica la forma ($\cup$ o $\cap$), el vèrtex, l'eix de simetria, els punts de tall amb els eixos i el recorregut de:\begin{tasks}(4)
  \task $y=x^2-6x+5$
  \task $y=-x^2+4x$
  \task $y=2x^2+4x+5$
  \task $y=-x^2-2x-1$
\end{tasks}\recorda{$V_x=\dfrac{-b}{2a}$, \ $V_y=f(V_x)$.}

\ex{1} Sense fer càlculs, digues quines d'aquestes paràboles tenen les branques cap amunt i quines cap avall. Quina és la més tancada? Quina té el vèrtex en $(0,\,1)$?\begin{tasks}(4)
  \task $y=3x^2$
  \task $y=-x^2$
  \task $y=\frac12x^2$
  \task $y=-2x^2+1$
\end{tasks}

\ex{2} Troba els punts de tall amb els eixos de les paràboles:\begin{tasks}(4)
  \task $y=x^2-x-6$
  \task $y=x^2-9$
  \task $y=2x^2-8x$
  \task $y=x^2+4x+4$
\end{tasks}

\ex{2} Comprova quins d'aquests punts són de la paràbola $y=x^2-3x+2$:\begin{tasks}(4)
  \task $(0,\,2)$
  \task $(1,\,0)$
  \task $(2,\,1)$
  \task $(-1,\,6)$
\end{tasks}

\ex{3} Troba l'expressió de la funció quadràtica en cada cas:\begin{tasks}(1)
  \task $y=x^2+bx+c$, sabent que passa pel punt $(0,\,3)$ i que té el vèrtex en $x=2$.
  \task $y=x^2+bx+c$, sabent que talla l'eix $OX$ en $x=3$ i en $x=-1$.
\end{tasks}


% =====================================================
\seccio{4}{Representació de funcions quadràtiques}

\exf{1} Representa les paràboles següents fent els càlculs necessaris:\begin{tasks}(3)
  \task $f(x)=x^2$
  \task $f(x)=x^2+1$
  \task $f(x)=x^2-2x-3$
  \task $f(x)=-x^2+4x-3$
  \task $f(x)=x^2+2x$
  \task $f(x)=-x^2+3$
\end{tasks}\recorda{forma ($\cup$ si $a>0$, $\cap$ si $a<0$), vèrtex i eix, tall amb $OY$ i el seu simètric, talls amb $OX$ i, si cal, una taula de valors al voltant del vèrtex.}

\ex{2} Representa $f(x)=-x^2+2x+3$ i indica'n el domini, el recorregut, on creix i on decreix, i el màxim o el mínim.

\ex{2} Representa $f(x)=x^2-4x+5$. Talla l'eix $OX$? Completa la representació amb una taula de valors.

\ex{2} Relaciona les gràfiques de paràboles següents amb les funcions corresponents, justificant la resposta:

\begin{center}\begin{tikzpicture}[x=0.5cm,y=0.5cm,baseline={(current bounding box.center)},font=\scriptsize]
\draw[muted!25,line width=.3pt] (-6,-6)--(-6,6);
\draw[muted!25,line width=.3pt] (-5,-6)--(-5,6);
\draw[muted!25,line width=.3pt] (-4,-6)--(-4,6);
\draw[muted!25,line width=.3pt] (-3,-6)--(-3,6);
\draw[muted!25,line width=.3pt] (-2,-6)--(-2,6);
\draw[muted!25,line width=.3pt] (-1,-6)--(-1,6);
\draw[muted!25,line width=.3pt] (0,-6)--(0,6);
\draw[muted!25,line width=.3pt] (1,-6)--(1,6);
\draw[muted!25,line width=.3pt] (2,-6)--(2,6);
\draw[muted!25,line width=.3pt] (3,-6)--(3,6);
\draw[muted!25,line width=.3pt] (4,-6)--(4,6);
\draw[muted!25,line width=.3pt] (5,-6)--(5,6);
\draw[muted!25,line width=.3pt] (6,-6)--(6,6);
\draw[muted!25,line width=.3pt] (7,-6)--(7,6);
\draw[muted!25,line width=.3pt] (8,-6)--(8,6);
\draw[muted!25,line width=.3pt] (-6,-6)--(8,-6);
\draw[muted!25,line width=.3pt] (-6,-5)--(8,-5);
\draw[muted!25,line width=.3pt] (-6,-4)--(8,-4);
\draw[muted!25,line width=.3pt] (-6,-3)--(8,-3);
\draw[muted!25,line width=.3pt] (-6,-2)--(8,-2);
\draw[muted!25,line width=.3pt] (-6,-1)--(8,-1);
\draw[muted!25,line width=.3pt] (-6,0)--(8,0);
\draw[muted!25,line width=.3pt] (-6,1)--(8,1);
\draw[muted!25,line width=.3pt] (-6,2)--(8,2);
\draw[muted!25,line width=.3pt] (-6,3)--(8,3);
\draw[muted!25,line width=.3pt] (-6,4)--(8,4);
\draw[muted!25,line width=.3pt] (-6,5)--(8,5);
\draw[muted!25,line width=.3pt] (-6,6)--(8,6);
\draw[-latex,muted] (-6,0)--(8,0);
\draw[-latex,muted] (0,-6)--(0,6);
\node[below=-1pt,text=muted] at (-6,0) {$-6$};
\node[below=-1pt,text=muted] at (-5,0) {$-5$};
\node[below=-1pt,text=muted] at (-4,0) {$-4$};
\node[below=-1pt,text=muted] at (-3,0) {$-3$};
\node[below=-1pt,text=muted] at (-2,0) {$-2$};
\node[below=-1pt,text=muted] at (-1,0) {$-1$};
\node[below=-1pt,text=muted] at (1,0) {$1$};
\node[below=-1pt,text=muted] at (2,0) {$2$};
\node[below=-1pt,text=muted] at (3,0) {$3$};
\node[below=-1pt,text=muted] at (4,0) {$4$};
\node[below=-1pt,text=muted] at (5,0) {$5$};
\node[below=-1pt,text=muted] at (6,0) {$6$};
\node[below=-1pt,text=muted] at (7,0) {$7$};
\node[below=-1pt,text=muted] at (8,0) {$8$};
\node[left=-1pt,text=muted] at (0,-6) {$-6$};
\node[left=-1pt,text=muted] at (0,-5) {$-5$};
\node[left=-1pt,text=muted] at (0,-4) {$-4$};
\node[left=-1pt,text=muted] at (0,-3) {$-3$};
\node[left=-1pt,text=muted] at (0,-2) {$-2$};
\node[left=-1pt,text=muted] at (0,-1) {$-1$};
\node[left=-1pt,text=muted] at (0,1) {$1$};
\node[left=-1pt,text=muted] at (0,2) {$2$};
\node[left=-1pt,text=muted] at (0,3) {$3$};
\node[left=-1pt,text=muted] at (0,4) {$4$};
\node[left=-1pt,text=muted] at (0,5) {$5$};
\node[left=-1pt,text=muted] at (0,6) {$6$};
\begin{scope}\clip (-6,-6) rectangle (8,6);
\draw[muted,line width=1.1pt,line join=round] (-6,-14)--(-5.888,-13.117)--(-5.776,-12.258)--(-5.664,-11.425)--(-5.552,-10.617)--(-5.44,-9.834)--(-5.328,-9.076)--(-5.216,-8.343)--(-5.104,-7.635)--(-4.992,-6.952)--(-4.88,-6.294)--(-4.768,-5.662)--(-4.656,-5.054)--(-4.544,-4.472)--(-4.432,-3.915)--(-4.32,-3.382)--(-4.208,-2.875)--(-4.096,-2.393)--(-3.984,-1.936)--(-3.872,-1.504)--(-3.76,-1.098)--(-3.648,-0.716)--(-3.536,-0.359)--(-3.424,-0.028)--(-3.312,0.279)--(-3.2,0.56)--(-3.088,0.816)--(-2.976,1.047)--(-2.864,1.254)--(-2.752,1.434)--(-2.64,1.59)--(-2.528,1.721)--(-2.416,1.827)--(-2.304,1.908)--(-2.192,1.963)--(-2.08,1.994)--(-1.968,1.999)--(-1.856,1.979)--(-1.744,1.934)--(-1.632,1.865)--(-1.52,1.77)--(-1.408,1.65)--(-1.296,1.504)--(-1.184,1.334)--(-1.072,1.139)--(-0.96,0.918)--(-0.848,0.673)--(-0.736,0.402)--(-0.624,0.107)--(-0.512,-0.214)--(-0.4,-0.56)--(-0.288,-0.931)--(-0.176,-1.327)--(-0.064,-1.748)--(0.048,-2.194)--(0.16,-2.666)--(0.272,-3.162)--(0.384,-3.683)--(0.496,-4.23)--(0.608,-4.802)--(0.72,-5.398)--(0.832,-6.02)--(0.944,-6.667)--(1.056,-7.339)--(1.168,-8.036)--(1.28,-8.758)--(1.392,-9.506)--(1.504,-10.278)--(1.616,-11.075)--(1.728,-11.898)--(1.84,-12.746)--(1.952,-13.618)--(2.064,-14.516)--(2.176,-15.439)--(2.288,-16.387)--(2.4,-17.36)--(2.512,-18)--(2.624,-18)--(2.736,-18)--(2.848,-18)--(2.96,-18)--(3.072,-18)--(3.184,-18)--(3.296,-18)--(3.408,-18)--(3.52,-18)--(3.632,-18)--(3.744,-18)--(3.856,-18)--(3.968,-18)--(4.08,-18)--(4.192,-18)--(4.304,-18)--(4.416,-18)--(4.528,-18)--(4.64,-18)--(4.752,-18)--(4.864,-18)--(4.976,-18)--(5.088,-18)--(5.2,-18)--(5.312,-18)--(5.424,-18)--(5.536,-18)--(5.648,-18)--(5.76,-18)--(5.872,-18)--(5.984,-18)--(6.096,-18)--(6.208,-18)--(6.32,-18)--(6.432,-18)--(6.544,-18)--(6.656,-18)--(6.768,-18)--(6.88,-18)--(6.992,-18)--(7.104,-18)--(7.216,-18)--(7.328,-18)--(7.44,-18)--(7.552,-18)--(7.664,-18)--(7.776,-18)--(7.888,-18)--(8,-18);
\draw[violet,line width=1.1pt,line join=round] (-6,18)--(-5.888,18)--(-5.776,18)--(-5.664,18)--(-5.552,18)--(-5.44,18)--(-5.328,18)--(-5.216,18)--(-5.104,18)--(-4.992,18)--(-4.88,18)--(-4.768,18)--(-4.656,18)--(-4.544,18)--(-4.432,18)--(-4.32,18)--(-4.208,18)--(-4.096,18)--(-3.984,17.856)--(-3.872,16.864)--(-3.76,15.898)--(-3.648,14.956)--(-3.536,14.039)--(-3.424,13.148)--(-3.312,12.281)--(-3.2,11.44)--(-3.088,10.624)--(-2.976,9.833)--(-2.864,9.066)--(-2.752,8.326)--(-2.64,7.61)--(-2.528,6.919)--(-2.416,6.253)--(-2.304,5.612)--(-2.192,4.997)--(-2.08,4.406)--(-1.968,3.841)--(-1.856,3.301)--(-1.744,2.786)--(-1.632,2.295)--(-1.52,1.83)--(-1.408,1.39)--(-1.296,0.976)--(-1.184,0.586)--(-1.072,0.221)--(-0.96,-0.118)--(-0.848,-0.433)--(-0.736,-0.722)--(-0.624,-0.987)--(-0.512,-1.226)--(-0.4,-1.44)--(-0.288,-1.629)--(-0.176,-1.793)--(-0.064,-1.932)--(0.048,-2.046)--(0.16,-2.134)--(0.272,-2.198)--(0.384,-2.237)--(0.496,-2.25)--(0.608,-2.238)--(0.72,-2.202)--(0.832,-2.14)--(0.944,-2.053)--(1.056,-1.941)--(1.168,-1.804)--(1.28,-1.642)--(1.392,-1.454)--(1.504,-1.242)--(1.616,-1.005)--(1.728,-0.742)--(1.84,-0.454)--(1.952,-0.142)--(2.064,0.196)--(2.176,0.559)--(2.288,0.947)--(2.4,1.36)--(2.512,1.798)--(2.624,2.261)--(2.736,2.75)--(2.848,3.263)--(2.96,3.802)--(3.072,4.365)--(3.184,4.954)--(3.296,5.568)--(3.408,6.206)--(3.52,6.87)--(3.632,7.559)--(3.744,8.274)--(3.856,9.013)--(3.968,9.777)--(4.08,10.566)--(4.192,11.381)--(4.304,12.22)--(4.416,13.085)--(4.528,13.975)--(4.64,14.89)--(4.752,15.83)--(4.864,16.794)--(4.976,17.785)--(5.088,18)--(5.2,18)--(5.312,18)--(5.424,18)--(5.536,18)--(5.648,18)--(5.76,18)--(5.872,18)--(5.984,18)--(6.096,18)--(6.208,18)--(6.32,18)--(6.432,18)--(6.544,18)--(6.656,18)--(6.768,18)--(6.88,18)--(6.992,18)--(7.104,18)--(7.216,18)--(7.328,18)--(7.44,18)--(7.552,18)--(7.664,18)--(7.776,18)--(7.888,18)--(8,18);
\draw[orange!85!black,line width=1.1pt,line join=round] (-6,-18)--(-5.888,-18)--(-5.776,-18)--(-5.664,-18)--(-5.552,-18)--(-5.44,-18)--(-5.328,-18)--(-5.216,-18)--(-5.104,-18)--(-4.992,-18)--(-4.88,-18)--(-4.768,-18)--(-4.656,-18)--(-4.544,-18)--(-4.432,-18)--(-4.32,-18)--(-4.208,-18)--(-4.096,-18)--(-3.984,-18)--(-3.872,-18)--(-3.76,-18)--(-3.648,-18)--(-3.536,-17.86)--(-3.424,-17.134)--(-3.312,-16.421)--(-3.2,-15.72)--(-3.088,-15.032)--(-2.976,-14.356)--(-2.864,-13.693)--(-2.752,-13.043)--(-2.64,-12.405)--(-2.528,-11.779)--(-2.416,-11.167)--(-2.304,-10.566)--(-2.192,-9.978)--(-2.08,-9.403)--(-1.968,-8.841)--(-1.856,-8.29)--(-1.744,-7.753)--(-1.632,-7.228)--(-1.52,-6.715)--(-1.408,-6.215)--(-1.296,-5.728)--(-1.184,-5.253)--(-1.072,-4.791)--(-0.96,-4.341)--(-0.848,-3.904)--(-0.736,-3.479)--(-0.624,-3.067)--(-0.512,-2.667)--(-0.4,-2.28)--(-0.288,-1.905)--(-0.176,-1.543)--(-0.064,-1.194)--(0.048,-0.857)--(0.16,-0.533)--(0.272,-0.221)--(0.384,0.078)--(0.496,0.365)--(0.608,0.639)--(0.72,0.901)--(0.832,1.15)--(0.944,1.386)--(1.056,1.61)--(1.168,1.822)--(1.28,2.021)--(1.392,2.207)--(1.504,2.381)--(1.616,2.542)--(1.728,2.691)--(1.84,2.827)--(1.952,2.951)--(2.064,3.062)--(2.176,3.161)--(2.288,3.247)--(2.4,3.32)--(2.512,3.381)--(2.624,3.429)--(2.736,3.465)--(2.848,3.488)--(2.96,3.499)--(3.072,3.497)--(3.184,3.483)--(3.296,3.456)--(3.408,3.417)--(3.52,3.365)--(3.632,3.3)--(3.744,3.223)--(3.856,3.134)--(3.968,3.031)--(4.08,2.917)--(4.192,2.79)--(4.304,2.65)--(4.416,2.497)--(4.528,2.333)--(4.64,2.155)--(4.752,1.965)--(4.864,1.763)--(4.976,1.548)--(5.088,1.32)--(5.2,1.08)--(5.312,0.827)--(5.424,0.562)--(5.536,0.284)--(5.648,-0.006)--(5.76,-0.309)--(5.872,-0.624)--(5.984,-0.952)--(6.096,-1.293)--(6.208,-1.646)--(6.32,-2.011)--(6.432,-2.389)--(6.544,-2.78)--(6.656,-3.183)--(6.768,-3.599)--(6.88,-4.027)--(6.992,-4.468)--(7.104,-4.921)--(7.216,-5.387)--(7.328,-5.866)--(7.44,-6.357)--(7.552,-6.86)--(7.664,-7.376)--(7.776,-7.905)--(7.888,-8.446)--(8,-9);
\draw[green!55!black,line width=1.1pt,line join=round] (-6,18)--(-5.888,18)--(-5.776,18)--(-5.664,18)--(-5.552,18)--(-5.44,18)--(-5.328,18)--(-5.216,18)--(-5.104,18)--(-4.992,18)--(-4.88,18)--(-4.768,18)--(-4.656,18)--(-4.544,18)--(-4.432,18)--(-4.32,18)--(-4.208,18)--(-4.096,18)--(-3.984,18)--(-3.872,18)--(-3.76,18)--(-3.648,18)--(-3.536,18)--(-3.424,18)--(-3.312,18)--(-3.2,18)--(-3.088,18)--(-2.976,18)--(-2.864,18)--(-2.752,18)--(-2.64,18)--(-2.528,18)--(-2.416,18)--(-2.304,18)--(-2.192,18)--(-2.08,18)--(-1.968,18)--(-1.856,18)--(-1.744,18)--(-1.632,18)--(-1.52,18)--(-1.408,18)--(-1.296,18)--(-1.184,17.506)--(-1.072,16.581)--(-0.96,15.682)--(-0.848,14.807)--(-0.736,13.958)--(-0.624,13.133)--(-0.512,12.334)--(-0.4,11.56)--(-0.288,10.811)--(-0.176,10.087)--(-0.064,9.388)--(0.048,8.714)--(0.16,8.066)--(0.272,7.442)--(0.384,6.843)--(0.496,6.27)--(0.608,5.722)--(0.72,5.198)--(0.832,4.7)--(0.944,4.227)--(1.056,3.779)--(1.168,3.356)--(1.28,2.958)--(1.392,2.586)--(1.504,2.238)--(1.616,1.915)--(1.728,1.618)--(1.84,1.346)--(1.952,1.098)--(2.064,0.876)--(2.176,0.679)--(2.288,0.507)--(2.4,0.36)--(2.512,0.238)--(2.624,0.141)--(2.736,0.07)--(2.848,0.023)--(2.96,0.002)--(3.072,0.005)--(3.184,0.034)--(3.296,0.088)--(3.408,0.166)--(3.52,0.27)--(3.632,0.399)--(3.744,0.554)--(3.856,0.733)--(3.968,0.937)--(4.08,1.166)--(4.192,1.421)--(4.304,1.7)--(4.416,2.005)--(4.528,2.335)--(4.64,2.69)--(4.752,3.07)--(4.864,3.474)--(4.976,3.905)--(5.088,4.36)--(5.2,4.84)--(5.312,5.345)--(5.424,5.876)--(5.536,6.431)--(5.648,7.012)--(5.76,7.618)--(5.872,8.248)--(5.984,8.904)--(6.096,9.585)--(6.208,10.291)--(6.32,11.022)--(6.432,11.779)--(6.544,12.56)--(6.656,13.366)--(6.768,14.198)--(6.88,15.054)--(6.992,15.936)--(7.104,16.843)--(7.216,17.775)--(7.328,18)--(7.44,18)--(7.552,18)--(7.664,18)--(7.776,18)--(7.888,18)--(8,18);
\node[text=muted] at (-3.9,-2.2) {\normalsize$\boldsymbol{f}$};
\node[text=violet] at (-1.9,3.8) {\normalsize$\boldsymbol{g}$};
\node[text=orange!85!black] at (5.9,0.7) {\normalsize$\boldsymbol{h}$};
\node[text=green!55!black] at (5.3,3.8) {\normalsize$\boldsymbol{i}$};
\end{scope}\end{tikzpicture}\end{center}
\begin{tasks}(4)
  \task $y=(x-3)^2$
  \task $y=(x-2)\cdot(x+1)$
  \task $y=-x^2-4x-2$
  \task $y=-\frac12x^2+3x-1$
\end{tasks}

\ex{3} Troba els punts on es tallen la paràbola $y=x^2-2x-3$ i la recta $y=x-3$. Representa-les per a comprovar-ho.


% =====================================================
\seccio{5}{Funcions radicals}

\exf{1} Completa la taula de valors de $y=\sqrt x$ i representa la funció. Quins són el domini i el recorregut? És creixent o decreixent?\begin{center}\renewcommand{\arraystretch}{1.25}\begin{tabular}{|c|C{1cm}|C{1cm}|C{1cm}|C{1cm}|C{1cm}|}\hline $x$ & $0$ & $1$ & $4$ & $9$ & $16$\\\hline $y$ & & & & & \\\hline\end{tabular}\end{center}

\ex{1} Calcula el domini de les funcions:\begin{tasks}(4)
  \task $y=\sqrt{x-5}$
  \task $y=\sqrt{x+3}$
  \task $y=\sqrt{2x-6}$
  \task $y=\sqrt{4-x}$
\end{tasks}\recorda{el radicand d'una arrel quadrada no pot ser negatiu: ha de ser $\geq0$.}

\ex{2} Indica el punt on comença cada gràfica, el domini i el recorregut, i representa-les amb una taula de valors:\begin{tasks}(4)
  \task $y=\sqrt{x-2}$
  \task $y=\sqrt x+1$
  \task $y=\sqrt{x+3}-2$
  \task $y=-\sqrt{x-1}$
\end{tasks}

\ex{2} Escriu l'expressió de la funció que s'obté en desplaçar la gràfica de $y=\sqrt x$:\begin{tasks}(1)
  \task 4 unitats cap avall.
  \task 2 unitats cap a l'esquerra.
  \task 3 unitats cap a la dreta i 1 cap amunt.
\end{tasks}

\ex{2} Donada la funció $f(x)=\sqrt{x+4}$, calcula $f(0)$, $f(5)$, $f(-4)$ i $f(-5)$. Quin valor de $x$ té per imatge 5?


% =====================================================
\seccio{6}{Funcions de proporcionalitat inversa}

\exf{1} Les magnituds $x$ i $y$ són inversament proporcionals. Completa la taula, escriu la funció que les relaciona i representa-la per a $x>0$.\begin{center}\renewcommand{\arraystretch}{1.25}\begin{tabular}{|c|C{1cm}|C{1cm}|C{1cm}|C{1cm}|C{1cm}|C{1cm}|}\hline $x$ & $1$ & $2$ & $3$ & $4$ & $6$ & $12$\\\hline $y$ & & $6$ & & & & \\\hline\end{tabular}\end{center}

\ex{1} Representa les funcions $y=\dfrac4x$ i $y=-\dfrac2x$ fent una taula de valors. Indica'n el domini, el recorregut, les asímptotes i si són creixents o decreixents.

\ex{1} Digues si cada taula correspon a una proporcionalitat directa, a una proporcionalitat inversa o a cap de les dues, i escriu-ne la fórmula:\begin{tasks}(2)
  \task \begin{center}\renewcommand{\arraystretch}{1.25}\begin{tabular}{|c|C{1cm}|C{1cm}|C{1cm}|C{1cm}|}\hline $x$ & $1$ & $2$ & $3$ & $6$\\\hline $y$ & $12$ & $6$ & $4$ & $2$\\\hline\end{tabular}\end{center}
  \task \begin{center}\renewcommand{\arraystretch}{1.25}\begin{tabular}{|c|C{1cm}|C{1cm}|C{1cm}|C{1cm}|}\hline $x$ & $1$ & $2$ & $3$ & $4$\\\hline $y$ & $5$ & $10$ & $15$ & $20$\\\hline\end{tabular}\end{center}
  \task \begin{center}\renewcommand{\arraystretch}{1.25}\begin{tabular}{|c|C{1cm}|C{1cm}|C{1cm}|C{1cm}|}\hline $x$ & $1$ & $2$ & $4$ & $5$\\\hline $y$ & $3$ & $5$ & $9$ & $11$\\\hline\end{tabular}\end{center}
  \task \begin{center}\renewcommand{\arraystretch}{1.25}\begin{tabular}{|c|C{1cm}|C{1cm}|C{1cm}|C{1cm}|}\hline $x$ & $2$ & $4$ & $5$ & $10$\\\hline $y$ & $10$ & $5$ & $4$ & $2$\\\hline\end{tabular}\end{center}
\end{tasks}

\ex{2} La gràfica de la funció $f(x)=\dfrac kx$ passa pel punt $(2,\,-3)$. Troba $k$, calcula $f(-4)$ i troba el valor de $x$ per al qual $f(x)=12$.

\ex{2} Quines d'aquestes funcions són de proporcionalitat inversa? Per a les que ho siguen, indica si són creixents o decreixents a cada branca.\begin{tasks}(4)
  \task $y=\dfrac5x$
  \task $y=\dfrac x5$
  \task $x\cdot y=-8$
  \task $y=5-x$
\end{tasks}


% =====================================================
\seccio{7}{Translacions de la funció de proporcionalitat inversa}

\exf{1} Indica les asímptotes, el domini i el recorregut d'aquestes hipèrboles, digues si són creixents o decreixents i representa-les:\begin{tasks}(4)
  \task $y=\dfrac{2}{x-3}$
  \task $y=\dfrac1x+2$
  \task $y=\dfrac{-3}{x+1}$
  \task $y=\dfrac{4}{x-2}-1$
\end{tasks}\recorda{l'asímptota vertical és on s'anul·la el denominador i l'horitzontal és el nombre que se suma a la fracció.}

\ex{1} Explica com s'obté la gràfica de cada funció a partir de la de $y=\dfrac3x$:\begin{tasks}(4)
  \task $y=\dfrac3x+4$
  \task $y=\dfrac3{x-5}$
  \task $y=\dfrac3{x+2}$
  \task $y=\dfrac3{x+1}-3$
\end{tasks}

\ex{2} Escriu l'expressió de la funció que s'obté en desplaçar la gràfica de $y=\dfrac2x$:\begin{tasks}(1)
  \task 3 unitats cap a la dreta.
  \task 1 unitat cap avall.
  \task 4 unitats cap a l'esquerra i 2 cap amunt.
\end{tasks}

\ex{2} Representa $f(x)=\dfrac{2}{x-1}+1$: dibuixa primer les asímptotes i després fes una taula de valors a banda i banda de l'asímptota vertical. On talla els eixos?

\ex{3} La gràfica següent és una hipèrbola del tipus $y=\dfrac{k}{x-a}+b$. Troba'n les asímptotes i l'expressió.

\begin{center}\begin{tikzpicture}[x=0.44cm,y=0.44cm,baseline={(current bounding box.center)},font=\scriptsize]
\draw[muted!25,line width=.3pt] (-4,-2)--(-4,8);
\draw[muted!25,line width=.3pt] (-3,-2)--(-3,8);
\draw[muted!25,line width=.3pt] (-2,-2)--(-2,8);
\draw[muted!25,line width=.3pt] (-1,-2)--(-1,8);
\draw[muted!25,line width=.3pt] (0,-2)--(0,8);
\draw[muted!25,line width=.3pt] (1,-2)--(1,8);
\draw[muted!25,line width=.3pt] (2,-2)--(2,8);
\draw[muted!25,line width=.3pt] (3,-2)--(3,8);
\draw[muted!25,line width=.3pt] (4,-2)--(4,8);
\draw[muted!25,line width=.3pt] (5,-2)--(5,8);
\draw[muted!25,line width=.3pt] (6,-2)--(6,8);
\draw[muted!25,line width=.3pt] (-4,-2)--(6,-2);
\draw[muted!25,line width=.3pt] (-4,-1)--(6,-1);
\draw[muted!25,line width=.3pt] (-4,0)--(6,0);
\draw[muted!25,line width=.3pt] (-4,1)--(6,1);
\draw[muted!25,line width=.3pt] (-4,2)--(6,2);
\draw[muted!25,line width=.3pt] (-4,3)--(6,3);
\draw[muted!25,line width=.3pt] (-4,4)--(6,4);
\draw[muted!25,line width=.3pt] (-4,5)--(6,5);
\draw[muted!25,line width=.3pt] (-4,6)--(6,6);
\draw[muted!25,line width=.3pt] (-4,7)--(6,7);
\draw[muted!25,line width=.3pt] (-4,8)--(6,8);
\draw[-latex,muted] (-4,0)--(6,0);
\draw[-latex,muted] (0,-2)--(0,8);
\node[below=-1pt,text=muted] at (-4,0) {$-4$};
\node[below=-1pt,text=muted] at (-3,0) {$-3$};
\node[below=-1pt,text=muted] at (-2,0) {$-2$};
\node[below=-1pt,text=muted] at (-1,0) {$-1$};
\node[below=-1pt,text=muted] at (1,0) {$1$};
\node[below=-1pt,text=muted] at (2,0) {$2$};
\node[below=-1pt,text=muted] at (3,0) {$3$};
\node[below=-1pt,text=muted] at (4,0) {$4$};
\node[below=-1pt,text=muted] at (5,0) {$5$};
\node[below=-1pt,text=muted] at (6,0) {$6$};
\node[left=-1pt,text=muted] at (0,-2) {$-2$};
\node[left=-1pt,text=muted] at (0,-1) {$-1$};
\node[left=-1pt,text=muted] at (0,1) {$1$};
\node[left=-1pt,text=muted] at (0,2) {$2$};
\node[left=-1pt,text=muted] at (0,3) {$3$};
\node[left=-1pt,text=muted] at (0,4) {$4$};
\node[left=-1pt,text=muted] at (0,5) {$5$};
\node[left=-1pt,text=muted] at (0,6) {$6$};
\node[left=-1pt,text=muted] at (0,7) {$7$};
\node[left=-1pt,text=muted] at (0,8) {$8$};
\begin{scope}\clip (-4,-2) rectangle (6,8);
\draw[accent,line width=1.1pt,line join=round] (-4,2.6)--(-3.967,2.597)--(-3.933,2.595)--(-3.9,2.592)--(-3.867,2.589)--(-3.833,2.586)--(-3.8,2.583)--(-3.767,2.58)--(-3.733,2.577)--(-3.7,2.574)--(-3.667,2.571)--(-3.633,2.568)--(-3.6,2.565)--(-3.567,2.562)--(-3.533,2.559)--(-3.5,2.556)--(-3.467,2.552)--(-3.433,2.549)--(-3.4,2.545)--(-3.367,2.542)--(-3.333,2.538)--(-3.3,2.535)--(-3.267,2.531)--(-3.233,2.528)--(-3.2,2.524)--(-3.167,2.52)--(-3.133,2.516)--(-3.1,2.512)--(-3.067,2.508)--(-3.033,2.504)--(-3,2.5)--(-2.967,2.496)--(-2.933,2.492)--(-2.9,2.487)--(-2.867,2.483)--(-2.833,2.478)--(-2.8,2.474)--(-2.767,2.469)--(-2.733,2.464)--(-2.7,2.459)--(-2.667,2.455)--(-2.633,2.45)--(-2.6,2.444)--(-2.567,2.439)--(-2.533,2.434)--(-2.5,2.429)--(-2.467,2.423)--(-2.433,2.417)--(-2.4,2.412)--(-2.367,2.406)--(-2.333,2.4)--(-2.3,2.394)--(-2.267,2.388)--(-2.233,2.381)--(-2.2,2.375)--(-2.167,2.368)--(-2.133,2.362)--(-2.1,2.355)--(-2.067,2.348)--(-2.033,2.341)--(-2,2.333)--(-1.967,2.326)--(-1.933,2.318)--(-1.9,2.31)--(-1.867,2.302)--(-1.833,2.294)--(-1.8,2.286)--(-1.767,2.277)--(-1.733,2.268)--(-1.7,2.259)--(-1.667,2.25)--(-1.633,2.241)--(-1.6,2.231)--(-1.567,2.221)--(-1.533,2.211)--(-1.5,2.2)--(-1.467,2.189)--(-1.433,2.178)--(-1.4,2.167)--(-1.367,2.155)--(-1.333,2.143)--(-1.3,2.13)--(-1.267,2.118)--(-1.233,2.104)--(-1.2,2.091)--(-1.167,2.077)--(-1.133,2.062)--(-1.1,2.048)--(-1.067,2.032)--(-1.033,2.016)--(-1,2)--(-0.967,1.983)--(-0.933,1.966)--(-0.9,1.947)--(-0.867,1.929)--(-0.833,1.909)--(-0.8,1.889)--(-0.767,1.868)--(-0.733,1.846)--(-0.7,1.824)--(-0.667,1.8)--(-0.633,1.776)--(-0.6,1.75)--(-0.567,1.723)--(-0.533,1.696)--(-0.5,1.667)--(-0.467,1.636)--(-0.433,1.605)--(-0.4,1.571)--(-0.367,1.537)--(-0.333,1.5)--(-0.3,1.462)--(-0.267,1.421)--(-0.233,1.378)--(-0.2,1.333)--(-0.167,1.286)--(-0.133,1.235)--(-0.1,1.182)--(-0.067,1.125)--(-0.033,1.065)--(0,1)--(0.033,0.931)--(0.067,0.857)--(0.1,0.778)--(0.133,0.692)--(0.167,0.6)--(0.2,0.5)--(0.233,0.391)--(0.267,0.273)--(0.3,0.143)--(0.333,0)--(0.367,-0.158)--(0.4,-0.333)--(0.433,-0.529)--(0.467,-0.75)--(0.5,-1)--(0.533,-1.286)--(0.567,-1.615)--(0.6,-2)--(0.633,-2.455)--(0.667,-3)--(0.7,-3.667)--(0.733,-4.5)--(0.767,-5.571)--(0.8,-7)--(0.833,-9)--(0.867,-12)--(0.9,-12)--(0.933,-12)--(0.967,-12);
\draw[accent,line width=1.1pt,line join=round] (0.975,-12);
\draw[accent,line width=1.1pt,line join=round] (0.983,-12);
\draw[accent,line width=1.1pt,line join=round] (0.992,-12);
\draw[accent,line width=1.1pt,line join=round] (1.008,18);
\draw[accent,line width=1.1pt,line join=round] (1.017,18);
\draw[accent,line width=1.1pt,line join=round] (1.025,18);
\draw[accent,line width=1.1pt,line join=round] (1.033,18)--(1.067,18)--(1.1,18)--(1.133,18)--(1.167,15)--(1.2,13)--(1.233,11.571)--(1.267,10.5)--(1.3,9.667)--(1.333,9)--(1.367,8.455)--(1.4,8)--(1.433,7.615)--(1.467,7.286)--(1.5,7)--(1.533,6.75)--(1.567,6.529)--(1.6,6.333)--(1.633,6.158)--(1.667,6)--(1.7,5.857)--(1.733,5.727)--(1.767,5.609)--(1.8,5.5)--(1.833,5.4)--(1.867,5.308)--(1.9,5.222)--(1.933,5.143)--(1.967,5.069)--(2,5)--(2.033,4.935)--(2.067,4.875)--(2.1,4.818)--(2.133,4.765)--(2.167,4.714)--(2.2,4.667)--(2.233,4.622)--(2.267,4.579)--(2.3,4.538)--(2.333,4.5)--(2.367,4.463)--(2.4,4.429)--(2.433,4.395)--(2.467,4.364)--(2.5,4.333)--(2.533,4.304)--(2.567,4.277)--(2.6,4.25)--(2.633,4.224)--(2.667,4.2)--(2.7,4.176)--(2.733,4.154)--(2.767,4.132)--(2.8,4.111)--(2.833,4.091)--(2.867,4.071)--(2.9,4.053)--(2.933,4.034)--(2.967,4.017)--(3,4)--(3.033,3.984)--(3.067,3.968)--(3.1,3.952)--(3.133,3.938)--(3.167,3.923)--(3.2,3.909)--(3.233,3.896)--(3.267,3.882)--(3.3,3.87)--(3.333,3.857)--(3.367,3.845)--(3.4,3.833)--(3.433,3.822)--(3.467,3.811)--(3.5,3.8)--(3.533,3.789)--(3.567,3.779)--(3.6,3.769)--(3.633,3.759)--(3.667,3.75)--(3.7,3.741)--(3.733,3.732)--(3.767,3.723)--(3.8,3.714)--(3.833,3.706)--(3.867,3.698)--(3.9,3.69)--(3.933,3.682)--(3.967,3.674)--(4,3.667)--(4.033,3.659)--(4.067,3.652)--(4.1,3.645)--(4.133,3.638)--(4.167,3.632)--(4.2,3.625)--(4.233,3.619)--(4.267,3.612)--(4.3,3.606)--(4.333,3.6)--(4.367,3.594)--(4.4,3.588)--(4.433,3.583)--(4.467,3.577)--(4.5,3.571)--(4.533,3.566)--(4.567,3.561)--(4.6,3.556)--(4.633,3.55)--(4.667,3.545)--(4.7,3.541)--(4.733,3.536)--(4.767,3.531)--(4.8,3.526)--(4.833,3.522)--(4.867,3.517)--(4.9,3.513)--(4.933,3.508)--(4.967,3.504)--(5,3.5)--(5.033,3.496)--(5.067,3.492)--(5.1,3.488)--(5.133,3.484)--(5.167,3.48)--(5.2,3.476)--(5.233,3.472)--(5.267,3.469)--(5.3,3.465)--(5.333,3.462)--(5.367,3.458)--(5.4,3.455)--(5.433,3.451)--(5.467,3.448)--(5.5,3.444)--(5.533,3.441)--(5.567,3.438)--(5.6,3.435)--(5.633,3.432)--(5.667,3.429)--(5.7,3.426)--(5.733,3.423)--(5.767,3.42)--(5.8,3.417)--(5.833,3.414)--(5.867,3.411)--(5.9,3.408)--(5.933,3.405)--(5.967,3.403)--(6,3.4);
\draw[muted,dashed,line width=.5pt] (1,-2)--(1,8);
\draw[muted,dashed,line width=.5pt] (-4,3)--(6,3);
\filldraw[fill=blue!70!black,draw=blue!70!black,line width=.8pt] (2,5) circle (1.5pt);
\filldraw[fill=blue!70!black,draw=blue!70!black,line width=.8pt] (3,4) circle (1.5pt);
\filldraw[fill=blue!70!black,draw=blue!70!black,line width=.8pt] (0,1) circle (1.5pt);
\end{scope}\end{tikzpicture}\end{center}



% =====================================================
\seccio{8}{Funcions exponencials}

\exf{1} Completa la taula de valors i representa en uns mateixos eixos $y=3^x$ i $y=\left(\frac13\right)^x$. Indica'n el domini, el recorregut, l'asímptota i si són creixents o decreixents.\begin{center}\renewcommand{\arraystretch}{1.3}\begin{tabular}{|c|C{1cm}|C{1cm}|C{1cm}|C{1cm}|C{1cm}|}\hline $x$ & $-2$ & $-1$ & $0$ & $1$ & $2$\\\hline $3^x$ & & & & & \\\hline $\left(\frac13\right)^x$ & & & & & \\\hline\end{tabular}\end{center}

\ex{1} Sense dibuixar-les, digues si aquestes funcions exponencials són creixents o decreixents:\begin{tasks}(4)
  \task $y=5^x$
  \task $y=0{,}4^x$
  \task $y=\left(\frac32\right)^x$
  \task $y=\left(\frac14\right)^x$
\end{tasks}

\ex{1} Per quin nombre es multiplica cada any una quantitat que…\begin{tasks}(2)
  \task augmenta un 6 \%?
  \task disminueix un 25 \%?
  \task es duplica?
  \task es redueix a la meitat?
\end{tasks}

\ex{2} Donada la funció $f(x)=2^x$, calcula $f(5)$, $f(-3)$ i $f(0)$. Quin valor de $x$ té per imatge 64? I $\frac14$?

\ex{2} Troba el valor de la lletra en cada cas:\begin{tasks}(1)
  \task La funció $y=a^x$ passa pel punt $(2,\,9)$.
  \task La funció $y=a^x$ passa pel punt $(-1,\,4)$.
  \task La funció $y=C\cdot2^x$ passa pel punt $(3,\,40)$.
\end{tasks}


% =====================================================
\seccio{9}{Problemes amb funcions}

\exf{1} Indica quin tipus de funció descriu cada situació i escriu-ne l'expressió:\begin{tasks}(1)
  \task El preu d'entrar a una fira que cobra 5 € per l'entrada i 2 € per cada atracció ($x$ atraccions).
  \task El temps que es tarda a recórrer 300 km segons la velocitat $x$ (en km/h).
  \task L'àrea d'un quadrat de costat $x$.
  \task El costat d'un quadrat d'àrea $x$.
  \task El nombre de seguidors d'un compte que en té 100 i es duplica cada mes ($x$ mesos).
\end{tasks}

\ex{1} Un taxi cobra 3,50 € per la baixada de bandera i 1,20 € per cada quilòmetre. Escriu la funció que dona el preu segons els quilòmetres recorreguts. Quant costa un trajecte de 8 km? Quants quilòmetres s'han fet si s'han pagat 21,50 €?

\ex{2} El gimnàs A cobra 20 € de matrícula i 25 € cada mes. El gimnàs B no cobra matrícula, però la quota és de 35 € cada mes. Escriu la funció que dona el que s'ha pagat en total en cada gimnàs al cap de $x$ mesos. Quan s'ha pagat el mateix? Quin convé més per a un any sencer?

\ex{2} Llancem una pilota cap amunt des de terra. La seua altura (en metres) al cap de $t$ segons és $h(t)=-5t^2+30t$.\begin{enumerate}[label=\alph*),leftmargin=*,itemsep=2pt,topsep=2pt]\item Quina altura màxima assoleix i en quin instant?\item Quan torna a terra?\item En quins instants està a 25 m d'altura?\end{enumerate}

\ex{2} De tots els rectangles de perímetre 36 cm, quin té l'àrea més gran? (Si un costat mesura $x$, quant mesura l'altre?)

\ex{3} Amb 60 m de tanca volem tancar un hort rectangular aprofitant una paret (només cal tancar tres costats). Si els dos costats perpendiculars a la paret mesuren $x$ metres, escriu la funció que dona l'àrea de l'hort. Quines dimensions fan que l'àrea siga màxima? Quant val?

\ex{2} Sis pintors tarden 10 dies a pintar un edifici. Escriu la funció que dona els dies que tardarien $n$ pintors treballant al mateix ritme. Quant tardarien 4 pintors? Quants pintors calen per a acabar en 5 dies?

\ex{2} La velocitat (en m/s) amb què arriba a terra un objecte que cau des d'una altura de $h$ metres és, aproximadament, $v(h)=\sqrt{20h}$. Amb quina velocitat arriba si cau des de 5 m? I des de 20 m i 45 m? Des de quina altura ha caigut si arriba a 40 m/s?

\ex{2} Un cultiu té 200 bacteris i es duplica cada hora. Escriu la funció que dona el nombre de bacteris al cap de $t$ hores. Quants n'hi haurà al cap de 5 hores? Quant de temps ha de passar perquè n'hi haja 51 200?

\ex{2} Depositem 5 000 € en un compte que dona un 2 \% d'interés compost anual. Escriu la funció que dona el capital al cap de $t$ anys i calcula quant hi haurà al cap de 3 anys i de 10 anys.

\ex{3} Un cotxe que val 18 000 € perd cada any el 15 \% del valor que tenia l'any anterior. Escriu la funció que dona el valor del cotxe al cap de $t$ anys. Quant valdrà d'ací a 2 anys? Quants anys han de passar perquè valga menys de la meitat?

\ex{3} Una botiga ven 200 samarretes al mes a 12 € cadascuna. Ha comprovat que, per cada euro que rebaixa el preu, en ven 40 més. Si rebaixa el preu $x$ euros, escriu la funció que dona els ingressos mensuals. Quin preu fa màxims els ingressos? Quants diners ingressa aleshores?

\vfill
\begin{center}\color{muted}\small\faCheckCircle[regular]\ \ Ànim! Comprova els teus resultats amb les eines interactives dels apunts de la unitat.\end{center}

\end{document}
